\exnumber{01}
\extitle{UI I - Ex 2 - Codes couleur}
\exfiles{ex01/fichier Figma, nomme Name\_FirstName\_DeliverableName\_Date\_VersionNumber}
\exauthorize{Ceux nommés dans l'exercice}
\makeheaderfiles

\textbf{Instructions :}

Allez sur ces 7 sites web et créez les palettes de couleurs à partir d'eux (3 à 4 couleurs principales).

\textbf{Livrable :}

\begin{itemize}
\item Un fichier Figma comprenant des captures de la page principale, l'URL, les couleurs extraites et les codes couleur correspondants.
\item Déposez le fichier Figma sur Github, nommé ainsi : « Name\_FirstName\_DeliverableName\_Date\_VersionNumber ». Exemple : « Doe\_John\_MidFiPrototype\_05242024\_V1 ».
\item Ajoutez un titre dans le fichier.
\item Soignez l'aspect global du fichier. Il doit être clair, simple et lisible. Vous pouvez regarder canva.com, sans surcharger le dessin de détails.
\end{itemize}

\textbf{Pistes :}

\begin{itemize}
\item Vous pouvez utiliser le document de l'exercice précédent pour ajouter les codes couleur sur Figma.
\item Utilisez la pipette de couleur pour détecter les codes couleur !
\end{itemize}

\textbf{Ressources :}

\begin{itemize}
\item Débuter en design graphique \url{https://edu.gcfglobal.org/en/beginning-graphic-design/}
\item Gareth graphic studio \url{https://www.youtube.com/c/GarethDavidStudio}
\end{itemize}

\begin{itemize}
\item https://www.komoot.com/ \url{https://www.komoot.com/} (App)
\item https://www.airtable.com/ \url{https://www.airtable.com/}
\item https://www.decathlon.com/collections/kids-gift-ideas \url{https://www.decathlon.com/collections/kids-gift-ideas}
\item https://www.jrpass.com/ \url{https://www.jrpass.com/}
\item https://www.thesiff.com/ \url{https://www.thesiff.com/}
\item https://brave.com/es/ \url{https://brave.com/es/}
\item https://www.nhp.gov.in/ \url{https://www.nhp.gov.in/}
\end{itemize}

